\section*{Capítulo \ref{ch:g8:balanced-equations} --- Conservação da massa e equações balanceadas}

\begin{solution}{exo:g8:balanced-equations:1}
Numa \omterm{def:g7:chemical-reactions:reaction}{reação química}, a massa total dos produtos formados é igual à massa
total dos \omterm{def:g7:chemical-reactions:reactant}{reagentes} consumidos.
\end{solution}

\begin{solution}{exo:g8:balanced-equations:2}
O coeficiente é 3. Ele representa 3 \omterm{def:g7:atoms-and-molecules:atom}{átomos} de carbono e $3 \times 2 = 6$
\omterm{def:g7:atoms-and-molecules:atom}{átomos} de oxigênio.
\end{solution}

\begin{solution}{exo:g8:balanced-equations:3}
Sim: 1 C e 2 O de cada lado.
\end{solution}

\begin{solution}{exo:g8:balanced-equations:4}
\ce{2Mg + O2 -> 2MgO}: 2 Mg e 2 O de cada lado.
\end{solution}

\begin{solution}{exo:g8:balanced-equations:5}
Nitrogênio: 2 N à esquerda, então 2 \ce{NH3}; ficam 6 H à direita, então
3 \ce{H2}: \ce{N2 + 3H2 -> 2NH3}.
\end{solution}

\begin{solution}{exo:g8:balanced-equations:6}
Pela \omterm{def:g8:balanced-equations:conservation}{conservação da massa}: $44 - 12 = \qty{32}{g}$ de oxigênio.
\end{solution}

\begin{solution}{exo:g8:balanced-equations:7}
O gás formado escapa para a sala, e a balança não o pesa mais. Nenhuma
massa é destruída: a massa que falta está no \omterm{def:g4:air-a-mixture-of-gases:air}{ar} da sala.
\end{solution}

\begin{solution}{exo:g8:balanced-equations:8}
Carbono: 1 de cada lado. Hidrogênio: 4 à esquerda, então 2 \ce{H2O}.
Oxigênio: $2 + 2 = 4$ à direita, então 2 \ce{O2}:
\ce{CH4 + 2O2 -> CO2 + 2H2O}.
\end{solution}

\begin{solution}{exo:g8:balanced-equations:9}
Transformar \ce{H2O} em \ce{H2O2} muda a substância: \ce{H2O2} é o
peróxido de hidrogênio (a água oxigenada), não a água. Só os
coeficientes podem ser mudados: \ce{2H2 + O2 -> 2H2O}.
\end{solution}

\begin{solution}{exo:g8:balanced-equations:10}
$20 \times 5 = 100$ \omterm{def:g7:atoms-and-molecules:molecule}{moléculas} de oxigênio; $20 \times 4 = 80$ \omterm{def:g7:atoms-and-molecules:molecule}{moléculas}
de água.
\end{solution}

\begin{solution}{exo:g8:balanced-equations:11}
Com 2 \ce{C2H6}: 4 C, então 4 \ce{CO2}; 12 H, então 6 \ce{H2O}; oxigênio
à direita $8 + 6 = 14$, então 7 \ce{O2}:
\ce{2C2H6 + 7O2 -> 4CO2 + 6H2O}.
\end{solution}

\begin{solution}{exo:g8:balanced-equations:12}
Continua marcando \qty{850.0}{g}. A caixa está lacrada: os
\qty{0.4}{g} de cera que queimaram viraram dióxido de carbono e vapor de
água, que continuam dentro da caixa junto com o oxigênio que eles
tomaram. Nada entrou nem saiu.
\end{solution}

\begin{solution}{pb:g8:balanced-equations:1}
\textbf{1.} Do oxigênio do \omterm{def:g4:air-a-mixture-of-gases:air}{ar}, que se combina com o ferro: o óxido de
ferro pesa o ferro mais esse oxigênio.

\textbf{2.} Tudo fica dentro do frasco: o oxigênio consumido já estava
no frasco, e o óxido formado fica lá. A massa total não muda.

\textbf{3.} A \omterm{def:g8:balanced-equations:conservation}{conservação da massa}: numa reação, a massa total dos
produtos é igual à massa total dos \omterm{def:g7:chemical-reactions:reactant}{reagentes}.

\textbf{4.} \ce{4Fe + 3O2 -> 2Fe2O3}.

\textbf{5.} \ce{3Fe + 2O2 -> Fe3O4}.

\textbf{6.} À esquerda: 4 Fe, $3 \times 2 = 6$ O. À direita: $2 \times 2 = 4$ Fe,
$2 \times 3 = 6$ O. Balanceada.

\textbf{7.} 4 \omterm{def:g7:atoms-and-molecules:atom}{átomos} de ferro para 3 \omterm{def:g7:atoms-and-molecules:molecule}{moléculas} de oxigênio, então $400 \times
\frac{3}{4} = 300$ \omterm{def:g7:atoms-and-molecules:molecule}{moléculas} de oxigênio.

\textbf{8.} $0.28 \times \frac{3}{7} = \qty{0.12}{g}$ de oxigênio.

\textbf{9.} $0.28 + 0.12 = \qty{0.40}{g}$ de óxido de ferro.

\textbf{10.} Sim: são necessários \qty{0.12}{g}, e o frasco continha
cerca de \qty{0.5}{g}.

\textbf{11.} O oxigênio consumido saiu do \omterm{def:g4:air-a-mixture-of-gases:air}{ar} do frasco (agora ele está
dentro do óxido sólido), então o frasco contém menos gás do que antes e
o \omterm{def:g4:air-a-mixture-of-gases:air}{ar} da sala entra. A leitura da balança então aumenta, da massa do \omterm{def:g4:air-a-mixture-of-gases:air}{ar}
que entra.

\textbf{12.} Foram tirados \qty{0.12}{g} de oxigênio do \omterm{def:g4:air-a-mixture-of-gases:air}{ar} do frasco.
\end{solution}
